% Historical source (removed by author request): references/_OLD_Übungsaufgaben/Übungsblatt 3 - Vektorechnung/Blatt3.tex:91
\begin{ExerciseBlock}
\ExerciseHeading{05}{Ausdrücke mit Skalar- und Vektorprodukten}{vector_algebra_product_expressions}
Nehmen Sie an, dass die Vektoren $\vec{a}$ und $\vec{b}$ senkrecht aufeinander stehen und die Längen  $|\vec{a}| = 1$ und $|\vec{b}| = 2$ haben. Berechnen Sie
\begin{enumerate}[label=\alph*)]
	\item $\left( 5 \vec{a} + \vec{b} \right) \cdot \left( 3 \vec{a} - 2 \vec{b} \right)$
	\item $\left| \left( 5 \vec{a} + \vec{b} \right) \times \left( 3 \vec{a} - 2 \vec{b} \right) \right|$
\end{enumerate} 

\end{ExerciseBlock}
\UseExerciseSolution{vector_algebra_product_expressions}
